% problem-000032 2 氨晶体多硫离子与热化学
\section*{2. 氨晶体多硫离子与热化学}
\textit{下列为分问。}

\subsection*{2-1}
氨晶体中，氨分⼦中的每个H均参与⼀个氢键的形成。N原⼦邻接⼏个氢原⼦？1摩尔固态氨中有⼏摩尔氢键？氨晶体融化时固态氨是下沉还是漂浮在液氨的液⾯上？

\subsection*{2-2}
$\mathrm { P } _ { 4 } \mathrm { S } _ { 5 }$ 是个多⾯体分⼦，结构中的多边形虽⾮平⾯状，但仍符合欧拉定律，两种原⼦成键后价层均满⾜8电⼦，S的氧化数为−2。画出该分⼦的结构图（⽤元素符号表⽰原⼦）。

\subsection*{2-3}
⽔煤⽓转化反应 $\mathrm { \small { [ C O ( g ) + H _ { 2 } O ( g )  H _ { 2 } ( g ) + C O _ { 2 } ( g ) ] } }$ ]是⼀个重要的化⼯过程。已知如下键能(BE)数据:$\mathrm { B E } ( \mathrm { C } \equiv \mathrm { O } ) = 1 0 7 2$ kJ $\mathrm { m o l ^ { - 1 } }$ $\mathrm { B E ( O { - } H ) } = 4 6 3 \mathrm { k J } \mathrm { m o l ^ { - 1 } }$ $\mathrm { B E } ( \mathrm { C = O } ) = 7 9 9$ kJ $\mathrm { m o l ^ { - 1 } }$ $\mathrm { B E } ( \mathrm { H } - \mathrm { H } ) = 4 3 6$ kJ$\mathrm { m o l } ^ { - 1 } \circ$ 估算反应热。该反应低温还是⾼温有利？简述理由。

\subsection*{2-4}
硫粉和S2−反应可以⽣成多硫离 $\yen 0$ 。在10mL S2−溶液中加⼊0.080g硫粉，控制条件使硫粉完全反应，检测到溶液中最⼤聚合度的多硫离⼦是 ${ \mathrm { S } } _ { 3 } { } ^ { 2 \cdot }$ −且 ${ \mathrm S } _ { n } ^ { 2 - } ~ ( n = 1 , 2 , 3 , \cdots )$ 离⼦浓度之⽐符合等⽐数列1, $1 0 , \cdots , 1 0 ^ { n - 1 }$ 。若不考虑其他副反应，计算反应后溶液中S2−的浓度 $\dot { \boldsymbol { { c } } } _ { 1 }$ 和其起始浓度 $\overset { \vartriangle } { \boldsymbol { c } _ { 0 ^ { \circ } } }$
